\chapter*{Preface}
\addcontentsline{toc}{chapter}{Preface}
\markboth{Preface}{Preface}

This book is about two ways of solving the same problem.

The first part solves it in software. You decide for yourself where a frame begins, how long it
is, how it is protected against errors, and what should happen when a reply never arrives. Every
mechanism becomes your own, and every mechanism becomes code you have to write and maintain. That
is instructive in a way that is hard to replace: you understand what a length field is for only
once you have forgotten one.

The second part throws all of that away and lets the hardware do the work. CAN dates from 1986,
it sits in every car built since the mid-nineties, and it solves very nearly the problems you have
just solved yourself --- except in silicon, and better. What is left for the software is talking
to the controller.

The third part connects it to reality: a microcontroller, an FPGA, four wires between them, and a
bus a bus analyser can read.

\section*{What makes the course unusual}

The CAN controller you write a driver for does not exist yet. It is being designed in VHDL by a
parallel class, at the same time as you write the software against it.

Neither class can compile against the other's code. The only thing binding the two halves
together is two documents --- a register map and an SPI protocol --- and they are the project's
real requirements specification.

\begin{aside}
\note Working against a contract instead of against an implementation is not a classroom
exercise. It is how embedded development works when hardware and software are built in parallel,
which is the normal case and not the exception. Most of what is hard in this course is hard for
that reason.
\end{aside}

\section*{How the book is arranged}

One chapter per lecture, fifteen of them, in three parts:

\begin{booktable}{@{}lL@{}}
\thead{Chapters} & \thead{Part} \\ \midrule
1--3 & \textbf{A protocol of your own.} Frames, byte parsing, buses, routing and reliability. \\
4--12 & \textbf{The CAN driver.} The register map, the architecture, the interface, the stub, SPI,
  the transport seam, the driver, the testing and the hardware. \\
13--15 & \textbf{The hardware.} Bring-up, bus analysis, and revision for the written test. \\
\end{booktable}

Every chapter opens with what it contains, ends with a summary, and has its exercises after that.
\textbf{The solutions to every exercise that does not ask for code are in appendix~A} --- all the
reasoning, calculation and decoding exercises, and the parts of the coding exercises that are
reasoning. The code that the coding exercises ask for is not answered there: for that, the
supplied test suite is the answer, and the code is published in the course repository after the
lecture in question.

The practice exam with full answers is last, in the appendix.

\section*{Assumed knowledge}

The book assumes classes, inheritance, virtual functions, abstract base classes and the factory
pattern. It does not revise them. It also assumes basic familiarity with Git, a terminal and
\code{make}.

Nothing about CAN, about hardware registers or about testing is assumed.

\section*{Further reading}

One more book belongs to the course, and it is not needed to pass it: \textit{CAN --- bussen,
framen och kontrollern}, the course's own book about CAN as a protocol. It covers the bus, the
frame, bit stuffing, CRC-15, arbitration, bit timing and error handling, with exercises and
answers, and is to be regarded as the long version of this book's chapter~4. It lives in a
repository of its own, because it is shared with the hardware course and therefore mentions no
course by name.
